\chapter{Сравнение существующих решений}


\section{Критерии сравнения}

Ниже перечислены критерии, которые используемые для сравнения различных типов рекомендательных систем.

\textbf{Тип используемых данных.} Определяет, какие данные лежат в основе работы метода.

\textbf{Наличие проблемы холодного старта.} Критерий показывает, может ли метод работать в ситуации, когда отсутствуют
данные о новых пользователях или объектах.

\textbf{Использование явных или неявных данных.} Оценивается, может ли метод обрабатывать только явные данные (например, оценки),
неявные данные (например, клики и просмотры) или оба типа данных.

\textbf{Возможность интерпретации рекомендаций.} Оценивается возможность объяснить пользователю,
почему была предложена определенная рекомендация.

\textbf{Разнообразие рекомендаций.} Определяет способность метода предлагать разнообразные
или новые объекты, а не только популярные элементы.

\textbf{Масштабируемость.} Этот критерий описывает, способна ли система эффективно работать
без изменений при росте количества пользователей и объектов.

\textbf{Потребность в доменной экспертизе.} Данный критерий определяет, требуется
ли для работы метода наличие специфических знаний о предметной области.

В таблице~\ref{tab:recommender_methods} приведено сравнение рекомендательных
систем по вышеперечисленным критериям

\begin{sidewaystable}
    \centering
    \begin{threeparttable}
        \caption{Сравнение методов рекомендательных систем по критериям}
        \label{tab:recommender_methods}
        \begin{tabular}{|l|c|c|c|c|c|}
            \hline
            \textbf{Критерий} & \textbf{\makecell{User-based \\ CF}} & \textbf{\makecell{Item-based \\ CF}} & \textbf{Content-based} & \textbf{\makecell{Knowledge- \\ based}} & \textbf{\makecell{Hybrid \\ Systems}} \\ \hline
            \textbf{\makecell{Тип данных}} & \makecell{Оценки \\ пользователей} & \makecell{Оценки \\ пользователей} & \makecell{Атрибуты \\ объектов}   & \makecell{Знания и      \\ примеры} & \makecell{Разнородные \\ данные} \\ \hline
            \textbf{\makecell{Проблема \\ холодного старта}} & Есть & Есть & Есть & Нет & Нет \\ \hline
            \textbf{\makecell{Использование \\ данных}} & Оба & Оба & Оба & Явные & Оба \\ \hline
            \textbf{Возможность интерпретации} & Нет & Да & Да & Да & \makecell{Зависит от \\ системы} \\ \hline
            \textbf{\makecell{Разнообразие \\ рекомендаций}} & Да & Да & Нет & Нет & Да \\ \hline
            \textbf{Масштабируемость} & Нет & Нет & Да & Да & Да \\ \hline
            \textbf{\makecell{Потребность в \\ экспертизе}} & Нет & Нет & Да & Да & \makecell{Зависит от \\ системы} \\ \hline
        \end{tabular}
    \end{threeparttable}
\end{sidewaystable}

\clearpage

\section{Результат сравнения}

Коллаборативная фильтрация не способна работать при отсутствии данных о
пользователях или объектах, так как она полностью зависит от истории взаимодействий.
Контентно-ориентированные методы способны справляться с холодным стартом, используя атрибуты
новых объектов. Методы, основанные на знаниях, также не зависят от истории взаимодействий и
могут рекомендовать товары на основе правил и ограничений. Гибридные системы, комбинируя методы,
эффективно решают проблему холодного старта.

Методы, основанные на знаниях, требуют наличия экспертной информации и правил для работы.
Контентно-ориентированные системы также требуют знания атрибутов объектов и их значимости для
пользователей. В коллаборативной фильтрации экспертные знания не требуются, так как метод
опирается на взаимодействия пользователей.

Методы, основанные на знаниях и контентно-ориентированные системы, легко интерпретируются,
так как объяснение рекомендаций можно связать с атрибутами или правилами.
Коллаборативная фильтрация менее интерпретируема, поскольку рекомендации основываются на
сходстве пользователей или объектов.

Таким образом, на основе проведенного сравнения можно сделать вывод, что
гибридные рекомендательные системы являются наиболее универсальными и гибкими,
так как они способны эффективно комбинировать данные различных типов,
решать проблему холодного старта и обрабатывать разреженные данные.
Однако их применение требует значительных вычислительных ресурсов и сложной архитектуры.
Методы коллаборативной фильтрации эффективны при наличии данных о
взаимодействиях, а контентно-ориентированные и основанные на знаниях системы полезны при
недостатке пользовательских данных и необходимости интерпретации рекомендаций.